% Options for packages loaded elsewhere
\PassOptionsToPackage{unicode}{hyperref}
\PassOptionsToPackage{hyphens}{url}
%
\documentclass[]{article}
\usepackage{amsmath,amssymb}
\usepackage{iftex}
\ifPDFTeX
  \usepackage[T1]{fontenc}
  \usepackage[utf8]{inputenc}
  \usepackage{textcomp}
\else
  \usepackage{unicode-math}
  \defaultfontfeatures{Scale=MatchLowercase}
  \defaultfontfeatures[\rmfamily]{Ligatures=TeX,Scale=1}
\fi
\usepackage{lmodern}
\IfFileExists{upquote.sty}{\usepackage{upquote}}{}
\IfFileExists{microtype.sty}{\usepackage[]{microtype}\UseMicrotypeSet[protrusion]{basicmath}}{}
\makeatletter
\@ifundefined{KOMAClassName}{%
  \IfFileExists{parskip.sty}{\usepackage{parskip}}{%
    \setlength{\parindent}{0pt}\setlength{\parskip}{6pt plus 2pt minus 1pt}}
}{\KOMAoptions{parskip=half}}
\makeatother
\usepackage{xcolor}
\usepackage{longtable,booktabs,array}
\usepackage{calc}
\usepackage{etoolbox}
\makeatletter
\patchcmd\longtable{\par}{\if@noskipsec\mbox{}\fi\par}{}{}
\makeatother
\IfFileExists{footnotehyper.sty}{\usepackage{footnotehyper}}{\usepackage{footnote}}
\makesavenoteenv{longtable}
\usepackage{graphicx}
\makeatletter
\def\maxwidth{\ifdim\Gin@nat@width>\linewidth\linewidth\else\Gin@nat@width\fi}
\def\maxheight{\ifdim\Gin@nat@height>\textheight\textheight\else\Gin@nat@height\fi}
\makeatother
\setkeys{Gin}{width=\maxwidth,height=\maxheight,keepaspectratio}
\makeatletter
\def\fps@figure{htbp}
\makeatother
\setlength{\emergencystretch}{3em}
\providecommand{\tightlist}{\setlength{\itemsep}{0pt}\setlength{\parskip}{0pt}}
\setcounter{secnumdepth}{-\maxdimen}
\ifLuaTeX\usepackage{selnolig}\fi
\IfFileExists{bookmark.sty}{\usepackage{bookmark}}{\usepackage{hyperref}}
\IfFileExists{xurl.sty}{\usepackage{xurl}}{}
\urlstyle{same}
\hypersetup{hidelinks,pdfcreator={LaTeX via pandoc}}

\title{}
\date{}

\begin{document}

\textbf{Aiko Memory System}

\textbf{Implementation, Evaluation, and Comparative Analysis}

\emph{A formulaic, auditable memory architecture for edge-deployed
conversational AI, benchmarked against Mem0, Zep/Graphiti, and
Letta/MemGPT}

Aiko Memory Research Series --- Paper II of II: Implementation \&
Evaluation

\textbf{Aiko-chan Development Team (AKA.\ OppaAI)}

Independent Research

August 2026 · Version 2.0

\emph{Companion volume: ``A Neuroscience-Grounded Computational Theory
of Personal AI Memory'' (Paper I, Version 2.0)}

\emph{Source of truth: the open-source Aiko-chan repository
(\texttt{cognition/memory/}) as of August 2026.}

\hypertarget{abstract}{%
\section{\texorpdfstring{\textbf{Abstract}}{Abstract}}\label{abstract}}

This paper documents the reference engineering implementation of the
memory theory developed in the companion volume (Paper I), and evaluates
it against both internal benchmarks and the current (2026) landscape of
agentic-memory frameworks --- Mem0, Zep/Graphiti, Letta (formerly
MemGPT), and related systems.

The system --- codenamed Aiko --- runs entirely on an 8\,GB-unified-memory
Jetson Orin Nano and rejects learned re-ranking in favor of closed-form,
auditable scoring. The live codebase implements:

\begin{itemize}
\item Ebbinghaus-style exponential decay with log-compressed access strength,
  valence-intensity modulation of effective half-life, spaced-repetition
  stretch from distinct access days, and optional mood-congruent
  modulation;
\item entity-centrality importance \(I_e=(1-\alpha)c(e)+\alpha r(e)\);
\item Reciprocal Rank Fusion (RRF) across KNN, FTS5 and entity-graph
  signals, plus pinned, recency, access, entity-importance, context-match
  and negative-penalty terms;
\item write-time valence/salience tagging and supersession-based belief
  revision;
\item a nightly dream pass (salience boost, near-duplicate merge, decay
  prune) and monthly consolidation with an LLM merge of survivors;
\item hierarchical layers including persona cache (identity facts) and
  scene summaries.
\end{itemize}

We document the SQLite + sqlite-vec + FTS5 schema, the full
ADD/RECALL/TOUCH/dream/consolidation pipeline, the testing and monitoring
harness, and end-to-end latency and footprint measurements drawn from a
year-long single-deployment field test.

On that deployment the system sustained \(\sim\)16\,ms median per-turn
retrieval latency, a \(\sim\)13\,MB total footprint after 1\,825 stored
facts, 0.72 precision@10 (vs.\ 0.68 for pure vector search), and 5.3\(\times\)
monthly archival compression with no manually-identified loss of
salience-tagged facts. We compare these figures, where methodologically
comparable, against publicly reported 2026 benchmark data for Mem0, Zep
and Letta, and discuss the engineering trade-offs --- auditability and
edge footprint versus leaderboard accuracy --- that motivate a formulaic
design.

\textbf{Keywords:} \emph{agentic memory, retrieval-augmented generation,
edge AI, SQLite, Reciprocal Rank Fusion, Mem0, Zep, Graphiti, Letta,
MemGPT, system evaluation}

\hypertarget{table-of-contents}{%
\section{\texorpdfstring{\textbf{Table of
Contents}}{Table of Contents}}\label{table-of-contents}}

\textbf{Abstract}

\textbf{1. Introduction \& Design Principles}

\textbf{2. System Architecture}

\textbf{3. The Core Formulas (Reference Implementation)}

\textbf{4. Data Schema}

\textbf{5. Worked End-to-End Example}

\textbf{6. Testing, Monitoring \& Tuning Harness}

\textbf{7. Performance Evaluation}

\textbf{8. Comparative Analysis vs.\ Established Agentic Memory Systems}

\textbf{9. Key Engineering Decisions \& Trade-offs}

\textbf{10. Limitations \& Future Work}

\textbf{11. Conclusion}

\textbf{References}

\textbf{Appendix A: Configuration Reference}

\hypertarget{introduction-design-principles}{%
\section{\texorpdfstring{\textbf{1. Introduction \& Design
Principles}}{1. Introduction \& Design Principles}}\label{introduction-design-principles}}

Paper I argued, from neuroscience and information theory, that
personal-AI memory should behave as a decaying, emotionally-modulated,
hierarchically-consolidated dynamical system rather than a static vector
index. This paper documents how that theory was turned into running code
in the open Aiko-chan repository, and reports how the resulting system
performs --- in absolute terms and relative to the current field.

\hypertarget{why-formulas-instead-of-learned-ranking}{%
\subsection{\texorpdfstring{\textbf{1.1 Why Formulas Instead of Learned
Ranking}}{1.1 Why Formulas Instead of Learned Ranking}}\label{why-formulas-instead-of-learned-ranking}}

The reference deployment target is a Jetson Orin Nano with 8\,GB of
unified memory and a 5--15\,W power envelope. Every additional resident
model competes with the primary conversational LLM for the same VRAM
budget; a learned-to-rank re-ranker would additionally require a labeled
training pipeline and periodic retraining. The system therefore uses
closed-form scoring whose principal terms are:

\begin{align*}
\text{strength} &= 1 + \log(1+\min(a, A_{\text{cap}})) \\
H_{\text{eff}} &= H_0 \cdot (1+\gamma\cdot I(v)) \cdot (1+\gamma_s\log(1+d)) \cdot m \\
s &= \text{strength} \cdot 0.5^{t/H_{\text{eff}}} \\
I(e) &= (1-\alpha)\,c(e) + \alpha\,r(e) \\
\text{RRF}(m) &= \sum_i \frac{1}{k+\text{rank}_i(m)} + \text{bonuses}
\end{align*}

Each formula is auditable (a developer can see exactly why a memory ranked
where it did), tunable by environment variables without retraining, and
cheap (microsecond-scale evaluation on the Jetson).

\hypertarget{design-constraints}{%
\subsection{\texorpdfstring{\textbf{1.2 Design
Constraints}}{1.2 Design Constraints}}\label{design-constraints}}

\begin{itemize}
\item Memory as lifecycle, not static storage --- facts age; high-intensity
  (especially negative) facts age more slowly; spacing and mood further
  modulate effective half-life.
\item Semantic + lexical + entity-graph fusion --- a memory should be
  findable via any one of three independent signals, fused by RRF.
\item Per-user isolation --- every table is keyed by \texttt{user\_id};
  optional SQLCipher encryption.
\item Efficient consolidation --- a nightly dream pass plus a monthly
  consolidation gate keep the store bounded.
\item Minimal context-window cost --- a typical recall is formatted to well
  under 1.2\,KB of injected context; a separate persona cache supplies
  stable identity facts every turn without a full search.
\end{itemize}

\hypertarget{system-architecture}{%
\section{\texorpdfstring{\textbf{2. System
Architecture}}{2. System Architecture}}\label{system-architecture}}

The live implementation lives under \texttt{cognition/memory/}. Hot-path
operations --- ADD, RECALL, TOUCH --- run on every conversational turn.
Offline maintenance comprises a nightly dream pass and a monthly
consolidation job. Supporting modules include \texttt{forget.py}
(decay), \texttt{entity.py} (importance, extraction, supersession,
valence/salience inference), \texttt{vecstore.py} (shared sqlite-vec +
RRF helpers), \texttt{schema.py} (DDL and tunables), episode/EMC
distillation, and optional cross-store / spreading / neural-net overlays.

\hypertarget{add-write-path}{%
\subsection{\texorpdfstring{\textbf{2.1 ADD (Write
Path)}}{2.1 ADD (Write Path)}}\label{add-write-path}}

\begin{itemize}
\item Extract candidate facts from the conversational turn (LLM-assisted
  extraction with identity-repair post-filters).
\item Embed each fact with the fixed 640-dimensional Harrier embedder
  (no fine-tuning).
\item Check near-duplicates via KNN against existing vectors (write-dedup
  threshold, default \(\approx 0.95\)).
\item Classify the write as add / no-op (exact duplicate) / supersede
  (belief revision).
\item Infer valence\_tag / valence\_score and salience\_hit with cheap
  heuristics (and optional LLM valence override).
\item Update the entity co-mention graph and enqueue the write
  asynchronously so it never blocks the turn.
\end{itemize}

\hypertarget{recall-read-path}{%
\subsection{\texorpdfstring{\textbf{2.2 RECALL (Read
Path)}}{2.2 RECALL (Read Path)}}\label{recall-read-path}}

\begin{itemize}
\item Run a tiered quick/wide pass across three signals: KNN (semantic),
  FTS5 (lexical, BM25-style), and entity graph (structural / spreading).
\item Fuse the rank lists via Reciprocal Rank Fusion (\(k=60\)).
\item Apply additive bonuses: pinned, recency, access, entity importance,
  mood-congruent context match; apply a soft negative-valence penalty
  unless the query itself is emotional/reflective.
\item Optionally re-rank the high-scoring subset by recency-among-relevant.
\item Touch each returned memory (increment \texttt{access\_count}, update
  \texttt{last\_accessed\_at}, update \texttt{access\_day\_count} for spacing).
\item Format the top results into a compact context block (bounded by
  \texttt{MEMORY\_CONTEXT\_TOTAL\_CHARS}, default 1200). A separate
  persona cache (kind=\texttt{identity}) is prepended every turn without
  a full search.
\end{itemize}

\hypertarget{maintenance-dream-pass-and-monthly-consolidation}{%
\subsection{\texorpdfstring{\textbf{2.3 Maintenance: Dream Pass and
Monthly
Consolidation}}{2.3 Maintenance: Dream Pass and Monthly Consolidation}}\label{maintenance-dream-pass-and-monthly-consolidation}}

The nightly dream pass (i) boosts \texttt{access\_count} on memories whose
composite salience score exceeds a threshold (stored \texttt{salience\_hit},
emotional intensity, access strength, spacing, mild recency prior, and
keyword policy), (ii) merges near-duplicate memories discovered via KNN
(retaining the higher-access / pinned copy and marking the other
superseded), and (iii) prunes memories whose decayed, access-weighted
score falls below the cleanup threshold after the grace period. Optional
EMC episode distillation converts undistrilled episodic traces into
durable semantic facts.

The monthly consolidation job applies a multi-factor retention gate
(salience, novelty, spacing, connectivity / entity importance) to the past
month's daily facts, LLM-merges the survivors into compact narrative
summaries, pins the result, and optionally retires the original daily
facts once merged.

\hypertarget{the-core-formulas-reference-implementation}{%
\section{\texorpdfstring{\textbf{3. The Core Formulas (Reference
Implementation)}}{3. The Core Formulas (Reference Implementation)}}\label{the-core-formulas-reference-implementation}}

\hypertarget{formula-1-ebbinghaus-decay-with-modulators}{%
\subsection{\texorpdfstring{\textbf{3.1 Formula 1 --- Ebbinghaus Decay
with Access, Emotion, Spacing and Mood
Modulators}}{3.1 Formula 1 --- Ebbinghaus Decay with Access, Emotion, Spacing and Mood Modulators}}\label{formula-1-ebbinghaus-decay-with-modulators}}

Implemented in \texttt{cognition/memory/forget.py}:

\begin{verbatim}
strength = 1 + log1p(min(access_count, ACCESS_COUNT_CAP))
H_eff = HALF_LIFE_DAYS
        * (1 + EMOTION_GAMMA * intensity(valence))
        * (1 + SPACING_GAMMA * log1p(access_day_count))
        * mood_factor
weighted_score = strength * (0.5 ** (days_since_last_access / H_eff))
\end{verbatim}

Intensity prefers \texttt{valence\_score} \(\in[-2,+2]\) mapped to
\(|\text{score}|/2\); falls back to tag intensities (default
neg\,=\,1.0, pos\,=\,0.4, neutral\,=\,0). Mood factor lengthens half-life
on valence match and shortens it on mismatch when a query valence is
supplied. Defaults: \(H_0=21\) days, cleanup threshold \(0.02\), grace
period 35 days, \(\gamma=0.5\).

Grace-protected memories (younger than \texttt{GRACE\_PERIOD\_DAYS}) are
never pruned.

\hypertarget{formula-2-entity-importance}{%
\subsection{\texorpdfstring{\textbf{3.2 Formula 2 --- Entity Importance
(Centrality +
Recency)}}{3.2 Formula 2 --- Entity Importance (Centrality + Recency)}}\label{formula-2-entity-importance}}

Implemented in \texttt{cognition/memory/entity.py}:

\[
I(e) = (1-\alpha)\,c(e) + \alpha\,r(e),\quad
c(e)=\frac{\text{comention}(e)}{\max_{e'} \text{comention}(e')},\quad
r(e)=e^{-\beta\cdot\Delta t(e)}
\]

with defaults \(\alpha=0.4\), \(\beta\approx0.05\) per day. The map is
cached briefly for recall and invalidated on writes. \(I_e\) participates
in the monthly retention gate and supplies a mild additive term in
recall ranking. Optional spreading activation and a neural-net overlay
(Hebbian edge growth, off by default) can further propagate activation
over the co-mention graph.

\hypertarget{formula-3-reciprocal-rank-fusion}{%
\subsection{\texorpdfstring{\textbf{3.3 Formula 3 --- Reciprocal Rank
Fusion (Multi-Signal
Fusion)}}{3.3 Formula 3 --- Reciprocal Rank Fusion (Multi-Signal Fusion)}}\label{formula-3-reciprocal-rank-fusion}}

\[
\text{RRF}(m)=\sum_{i\in\{\text{KNN,FTS,graph}\}}\frac{w_i}{k+\text{rank}_i(m)}
+\text{recency}+\text{access}+\text{pinned}+\text{entity imp.}
+\text{context match}-\text{neg penalty}
\]

with \(k=60\), graph weight default \(0.6\), and the remaining additive
weights exposed as environment variables. RRF is computed after a tiered
quick/wide candidate gather so that a memory that is only moderately
ranked on every signal can still surface.

\hypertarget{formula-4-emotional-imprinting-and-recall-modulation}{%
\subsection{\texorpdfstring{\textbf{3.4 Formula 4 --- Emotional Imprinting
and Recall
Modulation}}{3.4 Formula 4 --- Emotional Imprinting and Recall Modulation}}\label{formula-4-emotional-imprinting-and-recall-modulation}}

At write time, cheap lexical (and optional LLM) heuristics assign
\texttt{valence\_tag} \(\in\{\text{neg},\text{pos},\text{neutral}\}\) and
an integer \texttt{valence\_score} \(\in[-2,+2]\). These tags lengthen
\(H_{\text{eff}}\) proportionally to intensity and, at recall, apply a
soft proportional negative penalty unless the query itself carries
matching emotional or reflective intent. Positive intensity is present
but weaker than negative intensity by design (threat-asymmetric cost
structure, Paper I §4).

\hypertarget{formula-5-supersession}{%
\subsection{\texorpdfstring{\textbf{3.5 Formula 5 --- Supersession
(Belief
Revision)}}{3.5 Formula 5 --- Supersession (Belief Revision)}}\label{formula-5-supersession}}

Near-duplicate writes whose cosine similarity exceeds the dedup threshold
but whose normalized text differs are classified \texttt{supersede}. The
older row is marked \texttt{status='superseded'} (not deleted); the newer
row records \texttt{supersedes\_id}. Superseded facts are excluded from
ordinary recall but remain available for history / audit, including
optional chain expansion for identity/preference/plan kinds.

\hypertarget{formula-6-salience-and-consolidation}{%
\subsection{\texorpdfstring{\textbf{3.6 Formula 6 --- Salience Scoring and
Consolidation
Gate}}{3.6 Formula 6 --- Salience Scoring and Consolidation Gate}}\label{formula-6-salience-and-consolidation}}

Dream boost uses a composite \(0\ldots1\) salience score that combines
stored \texttt{salience\_hit}, emotional intensity, log-compressed access,
spacing days, mild recency, and a keyword-policy prior. Monthly
consolidation retains facts above a soft multi-factor threshold
(salience + novelty + spacing + connectivity) before LLM merge into pinned
narrative summaries.

\hypertarget{data-schema}{%
\section{\texorpdfstring{\textbf{4. Data
Schema}}{4. Data Schema}}\label{data-schema}}

\hypertarget{memories-core-table}{%
\subsection{\texorpdfstring{\textbf{4.1 memories (Core
Table)}}{4.1 memories (Core Table)}}\label{memories-core-table}}

\begin{longtable}[]{@{}
  >{\raggedright\arraybackslash}p{(\columnwidth - 4\tabcolsep) * \real{0.28}}
  >{\raggedright\arraybackslash}p{(\columnwidth - 4\tabcolsep) * \real{0.18}}
  >{\raggedright\arraybackslash}p{(\columnwidth - 4\tabcolsep) * \real{0.54}}@{}}
\toprule\noalign{}
\textbf{Column} & \textbf{Type} & \textbf{Notes} \\
\midrule\noalign{}
\endhead
\bottomrule\noalign{}
\endlastfoot
id & TEXT PK & UUID \\
user\_id & TEXT & Isolation boundary \\
memory & TEXT & Fact string \\
created\_at & TEXT & UTC ISO timestamp \\
access\_count & INTEGER & Touches, capped (default 255) \\
last\_accessed\_at & TEXT & UTC ISO, or \texttt{never} \\
pinned & INTEGER & 0/1; decay-proof when set \\
status & TEXT & \texttt{active} / \texttt{superseded} (default active) \\
supersedes\_id & TEXT & FK to superseded memory \\
kind & TEXT & \texttt{fact} / \texttt{identity} / \texttt{scene} /
\texttt{plan} / \texttt{event} / \texttt{schema} / \ldots \\
entities & TEXT & JSON array of entity strings \\
access\_day\_count & INTEGER & Distinct local days recalled (spacing) \\
valence\_tag & TEXT & \texttt{neg}/\texttt{pos}/\texttt{neutral} \\
valence\_score & INTEGER & \(-2\ldots+2\); NULL = legacy \\
salience\_hit & INTEGER & 0/1 keyword/policy flag at write \\
scene\_id & TEXT & L2 linkage to parent scene \\
\end{longtable}

\hypertarget{entity-relations-and-vectors}{%
\subsection{\texorpdfstring{\textbf{4.2 entity\_relations and
memories\_vec}}{4.2 entity\_relations and memories\_vec}}\label{entity-relations-and-vectors}}

\texttt{entity\_relations} stores co-mention edges
(\texttt{entity\_a}, \texttt{entity\_b}, \texttt{relation}, \texttt{weight},
\texttt{memory\_id}, \texttt{updated\_at}), scoped by \texttt{user\_id}.
\texttt{memories\_vec} stores L2-normalized 640-dimensional embeddings
keyed by memory id (sqlite-vec).

\hypertarget{worked-end-to-end-example}{%
\section{\texorpdfstring{\textbf{5. Worked End-to-End
Example}}{5. Worked End-to-End Example}}\label{worked-end-to-end-example}}

The following walkthrough traces one fact from ingestion through
consolidation, illustrating how the formulas interact in the current
codebase.

\textbf{Day 1, 10:00 --- ADD.} User: ``I'm starting a project called
GRACE --- a robotics assistant. I'll aim to finish it by Friday.'' Three
facts are extracted, identity-repaired, embedded, and written (no near
duplicate; all classified add). Entity edges GRACE\(\leftrightarrow\)robotics,
GRACE\(\leftrightarrow\)Friday, robotics\(\leftrightarrow\)Friday are recorded.
Valence tags are neutral; salience\_hit may fire on deadline-related
keywords.

\textbf{Day 1, 17:00 --- RECALL.} Query ``How's GRACE coming along?''
triggers quick/wide KNN, FTS and graph passes; RRF ranks the three
GRACE-related facts at the top. Each is touched (access\_count \(0\to1\),
last\_accessed\_at and access\_day\_count updated).

\textbf{Day 8 --- Decay check + new mention.} After seven days without
mention, strength \(\times 0.5^{7/H_{\text{eff}}}\) has declined, but grace
protection and any salience/access boosts keep the facts ranked highly.
A new fact (``Oppa completed GRACE ahead of the Friday deadline'') is
added (similarity below supersession threshold \(\Rightarrow\) add).

\textbf{Night of Day 8 --- Dream pass.} Composite salience elevates the
GRACE cluster; access\_count is boosted. Near-duplicate merge (if any)
retains the higher-access copy. Scores remain above the cleanup threshold,
so nothing is pruned.

\textbf{Month later --- Monthly consolidation.} The multi-factor gate
keeps the high-connectivity, deadline-bearing facts; LLM merge produces a
compact pinned monthly summary and the atomic daily rows may be retired.

\hypertarget{testing-monitoring-tuning-harness}{%
\section{\texorpdfstring{\textbf{6. Testing, Monitoring \& Tuning
Harness}}{6. Testing, Monitoring \& Tuning Harness}}\label{testing-monitoring-tuning-harness}}

Five production metrics gate health:

\begin{itemize}
\item \textbf{Retention rate} --- fraction of day-1-accessed memories still
  retrievable after 30 days (targets \(\ge70\%\) for salience-tagged,
  \(\ge50\%\) for routine).
\item \textbf{No premature forgetting} --- zero cleanups of facts younger
  than the grace period.
\item \textbf{Entity importance rank stability} --- top-10 entities drop
  at most two ranks over 30 days.
\item \textbf{RRF coverage} --- graph-only rescues appear in a non-trivial
  fraction of queries.
\item \textbf{Consolidation fidelity} --- zero critical facts lost;
  compression ratio in a healthy band.
\end{itemize}

Diagnostic playbook (symptoms \(\to\) likely cause \(\to\) fix) follows the
same pattern as v1: lengthen half-life or raise cleanup threshold if decay
is too aggressive; verify graph weight and entity extraction if the graph
signal is unused; re-run maintenance and inspect dedup threshold on
duplicate surfacing; raise consolidation soft threshold / max-month budget
if over-deletion occurs.

\hypertarget{performance-evaluation}{%
\section{\texorpdfstring{\textbf{7. Performance
Evaluation}}{7. Performance Evaluation}}\label{performance-evaluation}}

All figures are drawn from the same year-long, single-deployment field
test (Jetson Orin Nano, 8\,GB unified memory) documented in Paper I,
plus system-level instrumentation. As an \(N=1\) field deployment these
numbers characterize one system under one usage pattern.

\hypertarget{latency-budget}{%
\subsection{\texorpdfstring{\textbf{7.1 Latency
Budget}}{7.1 Latency Budget}}\label{latency-budget}}

Total per-turn context-injection latency --- persona cache (cached) +
KNN + FTS + graph + RRF scoring + touch metadata + formatting --- sums to
approximately 16\,ms median, comfortably inside a \(<100\)\,ms conversational
budget. Heavy operations (embedding calls, dream merge, monthly LLM merge)
remain off the hot path.

\hypertarget{storage-footprint}{%
\subsection{\texorpdfstring{\textbf{7.2 Storage
Footprint}}{7.2 Storage Footprint}}\label{storage-footprint}}

After \(\sim\)1\,825 facts the on-disk footprint of memories, 640-d
vectors, entity graph, FTS5 index and ancillary tables is approximately
13\,MB --- well under 1\,\% of the Jetson's 8\,GB envelope. The practical
constraint remains the LLM context window, not disk.

\hypertarget{retrieval-quality}{%
\subsection{\texorpdfstring{\textbf{7.3 Retrieval
Quality}}{7.3 Retrieval Quality}}\label{retrieval-quality}}

On a 100-query human-judged sample, RRF fusion improved precision@10 from
0.68 (pure KNN) to 0.72 and recall@10 from 0.61 to 0.68, while surfacing
an additional share of results that neither KNN nor FTS alone returned in
their top candidates (graph rescues).

\hypertarget{consolidation-efficiency}{%
\subsection{\texorpdfstring{\textbf{7.4 Consolidation
Efficiency}}{7.4 Consolidation Efficiency}}\label{consolidation-efficiency}}

Monthly compression (pre-LLM-merge survival \(\approx54.7\%\), post-merge
overall \(\approx5.3\times\)) with zero critical facts lost under manual
review of dropped rows. Persona cache remains a handful of stable identity
facts updated infrequently.

\hypertarget{comparative-analysis-vs.-established-agentic-memory-systems}{%
\section{\texorpdfstring{\textbf{8. Comparative Analysis vs.\ Established
Agentic Memory
Systems}}{8. Comparative Analysis vs.\ Established Agentic Memory Systems}}\label{comparative-analysis-vs.-established-agentic-memory-systems}}

By mid-2026 the agentic-memory space has consolidated around a handful of
widely adopted frameworks. Cross-system numbers should be read cautiously:
different projects report different harnesses, and third-party round-ups
are not a single controlled study.

\begin{longtable}[]{@{}
  >{\raggedright\arraybackslash}p{(\columnwidth - 6\tabcolsep) * \real{0.20}}
  >{\raggedright\arraybackslash}p{(\columnwidth - 6\tabcolsep) * \real{0.28}}
  >{\raggedright\arraybackslash}p{(\columnwidth - 6\tabcolsep) * \real{0.26}}
  >{\raggedright\arraybackslash}p{(\columnwidth - 6\tabcolsep) * \real{0.26}}@{}}
\toprule\noalign{}
\textbf{System} & \textbf{Core approach} & \textbf{Ranking} & \textbf{Typical footprint} \\
\midrule\noalign{}
\endhead
\bottomrule\noalign{}
\endlastfoot
Mem0 & Drop-in memory API; LLM extraction over a vector store & Embedding + heuristics & Hosted or self-hosted SDK \\
Zep / Graphiti & Temporal knowledge graph & Graph traversal + hybrid scoring & Graph DB; heavier self-hosting \\
Letta (MemGPT) & OS-inspired agent runtime (RAM / disk) & Agent-directed paging & Full agent runtime \\
Aiko (this work) & Formulaic decay/consolidation over SQLite + sqlite-vec + FTS5 & Closed-form RRF + modulators, no learned weights & \(\sim\)13\,MB/year; full edge on 8\,GB device \\
\end{longtable}

Aiko occupies the niche of hard edge budget + formula-level auditability
and emotional/psychological realism in retention. It does not claim to
match graph-reasoning depth (Zep) or self-managing runtime sophistication
(Letta) on shared benchmarks it has not yet run. The choice among these
systems remains an architectural commitment, not a checkbox.

\hypertarget{key-engineering-decisions-trade-offs}{%
\section{\texorpdfstring{\textbf{9. Key Engineering Decisions \&
Trade-offs}}{9. Key Engineering Decisions \& Trade-offs}}\label{key-engineering-decisions-trade-offs}}

\begin{itemize}
\item \textbf{No learned ranking.} Formulas stay inside the 8\,GB budget,
  need no training pipeline, and remain inspectable. Trade-off: possible
  sub-optimality on edge cases versus a well-trained re-ranker.
\item \textbf{Supersession instead of deletion.} Preserves belief history
  for audit without an LLM contradiction pass. Trade-off: modest extra
  store growth; superseded rows excluded from normal recall.
\item \textbf{Entity graph + optional spreading.} Rescues facts that are
  only findable through relationships. Trade-off: depends on reliable
  (currently rule-based) entity extraction.
\item \textbf{Separate monthly consolidation.} Nightly dream handles
  duplicates and ordinary decay; monthly compression handles volume.
  Trade-off: monthly facts are lossy by design.
\item \textbf{Log-compressed access strength.} Prevents heavily-touched
  memories from becoming effectively immortal while still rewarding
  spaced retrieval.
\end{itemize}

\hypertarget{limitations-future-work}{%
\section{\texorpdfstring{\textbf{10. Limitations \& Future
Work}}{10. Limitations \& Future Work}}\label{limitations-future-work}}

\begin{itemize}
\item Entity extraction remains largely rule-based; a lightweight NER or
  embedding-based linker would raise graph rescue rates.
\item Valence estimation is heuristic (with optional LLM override); sarcasm
  and implicit emotion remain imperfect.
\item Single-user validation (\(N=1\)); multi-user isolation is already
  implemented, but population-level claims require multi-subject study.
\item Decay parameters are global defaults; per-user Bayesian tuning is a
  natural extension.
\item Aiko has not yet been evaluated on LongMemEval or Deep Memory
  Retrieval; running those shared benchmarks is the highest-priority next
  step for quantitative comparison.
\item Optional neural-net overlay and full spreading activation are
  present but off or lightly weighted by default; their contribution under
  sustained multi-year load is still being measured.
\end{itemize}

\hypertarget{conclusion}{%
\section{\texorpdfstring{\textbf{11.
Conclusion}}{11. Conclusion}}\label{conclusion}}

Aiko's memory system turns the theoretical account of Paper I into a
working, fully auditable pipeline built on SQLite, sqlite-vec and FTS5,
running comfortably inside an 8\,GB edge-hardware budget. The closed-form
decay, importance, RRF, emotional, supersession and salience formulas
work together to give a system that is tunable without retraining,
explainable at every ranking decision, and lean enough to run indefinitely
on a single-board computer.

Measured against itself the system meets its design targets:
\(\sim\)16\,ms per-turn retrieval, \(\sim\)13\,MB/year storage growth, a
measurable precision gain from multi-signal fusion, and 5.3\(\times\) monthly
compression without loss of load-bearing facts. Measured against the
broader 2026 landscape it occupies a distinct niche complementary to
general-purpose, benchmark-optimized memory services.

\hypertarget{references}{%
\section{\texorpdfstring{\textbf{References}}{References}}\label{references}}

Cormack, G.\ V., Smucker, M.\ D., \& Clarke, C.\ L. (2009). Efficient and
effective spam filtering and re-ranking. Proceedings of SIGIR 2009,
713--714.

Robertson, S., \& Walker, S. (1994). Some simple effective approximations
to the 2-Poisson model. Proceedings of SIGIR 1994, 232--241.

Packer, C., et al. (2023). MemGPT: Towards LLMs as operating systems.
arXiv:2310.08560.

Rasmussen, P., et al. (2024--2025). Zep: A temporal knowledge graph
architecture for agent memory (Graphiti). Zep technical report.

Mem0 project documentation, mem0.ai (accessed 2026).

Letta (formerly MemGPT) project documentation, letta.com (accessed 2026).

Aiko-chan source repository, github.com/OppaAI/Aiko-chan (dev branch,
August 2026) --- especially \texttt{cognition/memory/\{forget,entity,schema,vecstore,episode\_dream\}.py}.

\hypertarget{appendix-a-configuration-reference}{%
\section{\texorpdfstring{\textbf{Appendix A: Configuration
Reference}}{Appendix A: Configuration Reference}}\label{appendix-a-configuration-reference}}

\begin{longtable}[]{@{}
  >{\raggedright\arraybackslash}p{(\columnwidth - 6\tabcolsep) * \real{0.40}}
  >{\raggedright\arraybackslash}p{(\columnwidth - 6\tabcolsep) * \real{0.14}}
  >{\raggedright\arraybackslash}p{(\columnwidth - 6\tabcolsep) * \real{0.16}}
  >{\raggedright\arraybackslash}p{(\columnwidth - 6\tabcolsep) * \real{0.30}}@{}}
\toprule\noalign{}
\textbf{Parameter} & \textbf{Default} & \textbf{Range / notes} & \textbf{Interpretation} \\
\midrule\noalign{}
\endhead
\bottomrule\noalign{}
\endlastfoot
FORGET\_HALF\_LIFE\_DAYS & 21.0 & [7, 60] & Base half-life (days) \\
FORGET\_CLEANUP\_THRESHOLD & 0.02 & [0.01, 0.1] & Score floor for prune \\
FORGET\_GRACE\_PERIOD\_DAYS & 35 & [14, 60] & Protect new facts \\
FORGET\_ACCESS\_COUNT\_CAP & 255 & --- & Cap for log1p strength \\
FORGET\_EMOTION\_GAMMA & 0.5 & [0, 1] & Valence stretch of \(H_{\text{eff}}\) \\
FORGET\_INTENSITY\_NEG / POS / NEUTRAL & 1.0 / 0.4 / 0.0 & --- & Tag intensities \\
FORGET\_MOOD\_MATCH\_SLOWDOWN & 1.3 & --- & Match lengthens \(H_{\text{eff}}\) \\
FORGET\_MOOD\_MISMATCH\_ACCELERATION & 1.3 & --- & Mismatch shortens \(H_{\text{eff}}\) \\
RRF\_K & 60 & [30, 100] & RRF damping constant \\
MEMORY\_RANK\_GRAPH\_WEIGHT & 0.6 & [0, 1] & Graph contribution \\
ENTITY\_IMPORTANCE\_ALPHA & 0.4 & [0.2, 0.6] & Recency vs centrality \\
MEMORY\_RANK\_ENTITY\_IMPORTANCE\_WEIGHT & 0.008 & --- & Mild recall boost \\
WRITE\_DEDUP / DREAM\_MERGE thresholds & \(\approx0.95\) / \(0.88\) & --- & Supersede / merge \\
CONSOLIDATION soft threshold / max month & \(\approx0.4\) / 30 & --- & Monthly gate \\
EMBED\_DIMS & 640 & --- & Harrier embedding dimension \\
\end{longtable}

\end{document}
